% GENERATED FILE. Edit canonical lessons, then run npm run generate:books.
% canonical-source-hash: fnv1a64:c16db459a7e41937
% canonical-lessons: FR-C516-pinceau, FR-C516-toile, FR-C516-aquarelle, FR-C516-chevalet, FR-C516-esquisser

\chapter{Pinceaux et couleurs}
\label{ch:fr-pinceaux-et-couleurs}
\hlchaptermodality{516}

\begin{chapteropening}
\textbf{By the end of this chapter:} \emph{I can talk about painting and drawing.}

Complete the last lesson of chapter 516: I can talk about painting and drawing.
\end{chapteropening}

\section[le pinceau]{le pinceau — paintbrush}

\label{lesson:FR-C516-pinceau}



\begin{quote}
Before the new one: say the French for actor, then the French for applause.
\end{quote}

\subsection*{You'll want to know: le pinceau}
\textbf{le pinceau} — \textquotedblleft{}paintbrush\textquotedblright{}.

\begin{quote}
L'artiste nettoie ses pinceaux après le travail.
\end{quote}

\begin{quote}
\emph{The artist cleans her brushes after work.}
\end{quote}

Plural \emph{les pinceaux}. For walls, a bigger brush is \emph{une brosse}.

\begin{cousinweb}
\emph{le pinceau} comes from Old French \emph{pincel} (EtymDB), from Vulgar Latin \emph{pēnicellus}, a form of Latin \emph{pēniculus} \textquotedblleft{}a brush\textquotedblright{} (Lewis \& Short).
\end{cousinweb}

\begin{grammarlens}[title={What you've built}]
One more word for telling, explaining and describing at B1.
\end{grammarlens}

\subsection*{Your turn}
\begin{itemize}
\raggedright
  \item \emph{Say it:} \emph{le pinceau}
  \item \emph{Say it:} \emph{le pinceau}, once more
  \item \emph{Say it:} a sentence of your own with \emph{le pinceau}
  \item \emph{Recall:} say the French for actor, then the French for applause, then say \emph{le pinceau} again
\end{itemize}

\begin{tcolorbox}[breakable,colback=teal!4,colframe=teal!35!black,title={Before you move on}]
What does le pinceau mean? (\textquotedblleft{}Paintbrush\textquotedblright{}.) Say it once more.
\end{tcolorbox}
\section[la toile]{la toile — canvas; cloth; web}

\label{lesson:FR-C516-toile}



\begin{quote}
Before the new one: say the French for applause, then the French for paintbrush.
\end{quote}

\subsection*{You'll want to know: la toile}
\textbf{la toile} — \textquotedblleft{}canvas; cloth; web\textquotedblright{}.

\begin{quote}
Le peintre a vendu sa plus grande toile.
\end{quote}

\begin{quote}
\emph{The painter sold his biggest canvas.}
\end{quote}

\emph{Une toile d'araignée} — a spider's web. \emph{La Toile} — the Web.

\begin{cousinweb}
\emph{la toile} is inherited from Latin \emph{tēla} \textquotedblleft{}a web, a cloth on the loom\textquotedblright{} (EtymDB), related to \emph{texere} \textquotedblleft{}to weave\textquotedblright{} (Lewis \& Short).
\end{cousinweb}

\begin{grammarlens}[title={What you've built}]
One more word for telling, explaining and describing at B1.
\end{grammarlens}

\subsection*{Your turn}
\begin{itemize}
\raggedright
  \item \emph{Say it:} \emph{la toile}
  \item \emph{Say it:} \emph{la toile}, once more
  \item \emph{Say it:} a sentence of your own with \emph{la toile}
  \item \emph{Recall:} say the French for applause, then the French for paintbrush, then say \emph{la toile} again
\end{itemize}

\begin{tcolorbox}[breakable,colback=teal!4,colframe=teal!35!black,title={Before you move on}]
What does la toile mean? (\textquotedblleft{}Canvas; cloth; web\textquotedblright{}.) Say it once more.
\end{tcolorbox}
\section[l'aquarelle]{l'aquarelle — watercolour}

\label{lesson:FR-C516-aquarelle}



\begin{quote}
Before the new one: say the French for paintbrush, then the French for canvas.
\end{quote}

\subsection*{You'll want to know: l'aquarelle}
\textbf{l'aquarelle} — \textquotedblleft{}watercolour\textquotedblright{}.

\begin{quote}
Elle peint des paysages à l'aquarelle.
\end{quote}

\begin{quote}
\emph{She paints landscapes in watercolour.}
\end{quote}

Feminine. Both the technique and a painting.

\begin{cousinweb}
\emph{l'aquarelle} was borrowed from Italian \emph{acquarella}, built on \emph{acqua} \textquotedblleft{}water\textquotedblright{}, inherited from Latin \emph{aqua} (Lewis \& Short).
\end{cousinweb}

\begin{grammarlens}[title={What you've built}]
One more word for telling, explaining and describing at B1.
\end{grammarlens}

\subsection*{Your turn}
\begin{itemize}
\raggedright
  \item \emph{Say it:} \emph{l'aquarelle}
  \item \emph{Say it:} \emph{l'aquarelle}, once more
  \item \emph{Say it:} a sentence of your own with \emph{l'aquarelle}
  \item \emph{Recall:} say the French for paintbrush, then the French for canvas, then say \emph{l'aquarelle} again
\end{itemize}

\begin{tcolorbox}[breakable,colback=teal!4,colframe=teal!35!black,title={Before you move on}]
What does l'aquarelle mean? (\textquotedblleft{}Watercolour\textquotedblright{}.) Say it once more.
\end{tcolorbox}
\section[le chevalet]{le chevalet — easel}

\label{lesson:FR-C516-chevalet}



\begin{quote}
Before the new one: say the French for canvas, then the French for watercolour.
\end{quote}

\subsection*{You'll want to know: le chevalet}
\textbf{le chevalet} — \textquotedblleft{}easel\textquotedblright{}.

\begin{quote}
Le peintre installe son chevalet au bord de la rivière.
\end{quote}

\begin{quote}
\emph{The painter sets up his easel by the river.}
\end{quote}

Literally a little horse.

\begin{cousinweb}
\emph{le chevalet} is built on \emph{cheval} \textquotedblleft{}horse\textquotedblright{} with the small ending \emph{-et} (EtymDB); \emph{cheval} is inherited from Latin \emph{caballus} \textquotedblleft{}a horse\textquotedblright{} (Lewis \& Short).
\end{cousinweb}

\begin{grammarlens}[title={What you've built}]
One more word for telling, explaining and describing at B1.
\end{grammarlens}

\subsection*{Your turn}
\begin{itemize}
\raggedright
  \item \emph{Say it:} \emph{le chevalet}
  \item \emph{Say it:} \emph{le chevalet}, once more
  \item \emph{Say it:} a sentence of your own with \emph{le chevalet}
  \item \emph{Recall:} say the French for canvas, then the French for watercolour, then say \emph{le chevalet} again
\end{itemize}

\begin{tcolorbox}[breakable,colback=teal!4,colframe=teal!35!black,title={Before you move on}]
What does le chevalet mean? (\textquotedblleft{}Easel\textquotedblright{}.) Say it once more.
\end{tcolorbox}
\section[esquisser]{esquisser — to sketch, to outline}

\label{lesson:FR-C516-esquisser}



\begin{quote}
Before the new one: say the French for watercolour, then the French for easel.
\end{quote}

\subsection*{You'll want to know: esquisser}
\textbf{esquisser} — \textquotedblleft{}to sketch, to outline\textquotedblright{}.

\begin{quote}
L'architecte a esquissé la maison en quelques minutes.
\end{quote}

\begin{quote}
\emph{The architect sketched the house in a few minutes.}
\end{quote}

\emph{Une esquisse} — a sketch. \emph{Esquisser un sourire} — to give a faint smile.

\begin{cousinweb}
\emph{esquisser} is built on \emph{esquisse} \textquotedblleft{}a sketch\textquotedblright{}, borrowed from Italian \emph{schizzo} (EtymDB), whose origin is usually traced to a Latin word for a sudden improvised poem.
\end{cousinweb}

\begin{grammarlens}[title={What you've built}]
One more word for telling, explaining and describing at B1.
\end{grammarlens}

\subsection*{Your turn}
\begin{itemize}
\raggedright
  \item \emph{Say it:} \emph{esquisser}
  \item \emph{Say it:} \emph{esquisser}, once more
  \item \emph{Say it:} a sentence of your own with \emph{esquisser}
  \item \emph{Recall:} say the French for watercolour, then the French for easel, then say \emph{esquisser} again
\end{itemize}

\begin{tcolorbox}[breakable,colback=teal!4,colframe=teal!35!black,title={Before you move on}]
What does esquisser mean? (\textquotedblleft{}To sketch, to outline\textquotedblright{}.) Say it once more.
\end{tcolorbox}
